% BHU PYQ -- Professional Exam Template
\documentclass[11pt,a4paper]{article}

\usepackage[a4paper,top=2.5cm,bottom=2.5cm,left=2.8cm,right=2.8cm,headheight=14pt]{geometry}
\usepackage{amsmath,amssymb}
\usepackage[T1]{fontenc}
\usepackage[utf8]{inputenc}
\usepackage{lmodern}
\usepackage{microtype}
\usepackage{fancyhdr}
\usepackage{enumitem}
\usepackage{hyperref}
\usepackage{lastpage}
\usepackage{parskip}

\hypersetup{colorlinks=true,urlcolor=black,linkcolor=black,
  pdftitle={B.Sc. (Hons.) Botany Sem-V BOB-502 -- BHU PYQ}}

\pagestyle{fancy}
\fancyhf{}
\renewcommand{\headrulewidth}{0.4pt}
\renewcommand{\footrulewidth}{0.4pt}
\fancyhead[L]{\small   \textit{Banaras Hindu University}}
\fancyhead[R]{\small   \textit{Botany Paper: BOB-502}}
\fancyfoot[C]{\small Page  \thepage\ of \pageref{LastPage}}


\newlist{parts}{enumerate}{2}
\setlist[parts,1]{label=(\alph*),leftmargin=*,itemsep=4pt,topsep=3pt}
\setlist[parts,2]{label=(\roman*),leftmargin=*,itemsep=2pt,topsep=2pt}

\newcommand{\pts}[1]{\hfill   \textit{[#1]}}


\begin{document}


\begin{center}
  {\large\bfseries BANARAS HINDU UNIVERSITY}\\[2pt]
  {
\normalsize B.Sc. (Hons.) Semester V Examination, 2016-17}\\[6pt]
  \rule{0.8\linewidth}{0.5pt}\\[6pt]
  {\large\bfseries Botany}\\[4pt]
  {\large\bfseries Paper: BOB-502 --- Comparative Studies of Phanerogams}\\[4pt]
  {
\normalsize Time: Three Hours \hfill Full Marks: 70}
\end{center}

\rule{\linewidth}{0.4pt}

\smallskip



\noindent   \textit{Note: Attempt five questions in all, selecting two questions from each Section-A and B. Question No. 9 (Section-C) is compulsory. The figures in the right-hand margin indicate marks.}

\bigskip

\begin{center}
   \textbf{SECTION -- A}
\end{center}

\medskip


\noindent   \textbf{Question 1.} Give a detailed account of the order Gnetales with special reference to   \textit{Gnetum}.\pts{15}

\medskip


\noindent   \textbf{Question 2.} Discuss briefly any two of the following:\pts{$7.5  \times 2 = 15$}

\begin{parts}
  \item Embryogeny in conifers
  \item   \textit{Pentoxylon}
  \item Origin of Gymnosperms
\end{parts}

\medskip


\noindent   \textbf{Question 3.} With the help of suitable diagrams, describe briefly the morphology and reproduction in   \textit{Zamia}.\pts{15}

\medskip


\noindent   \textbf{Question 4.} With the help of labelled diagrams, differentiate between any three of the following:\pts{$5  \times 3 = 15$}

\begin{parts}
  \item Male and female strobili of   \textit{Ginkgo}
  \item Anatomy of leaves of   \textit{Gnetum} and   \textit{Zamia}
  \item Manoxylic and pycnoxylic wood
  \item Female cones of   \textit{Zamia} and   \textit{Biota}
\end{parts}

\bigskip

\begin{center}
   \textbf{SECTION -- B}
\end{center}

\medskip


\noindent   \textbf{Question 5.} With the help of floral formula and floral diagram, describe the family Scrophulariaceae. Give names of two economically important plants of this family.\pts{15}

\medskip


\noindent   \textbf{Question 6.} Write short notes on any two of the following:\pts{$7.5  \times 2 = 15$}

\begin{parts}
  \item Anther culture
  \item Outline of Hutchinson's system of classification
  \item Numerical taxonomy
\end{parts}

\medskip


\noindent   \textbf{Question 7.} Write salient features of the family Rutaceae and discuss its economic importance.\pts{15}

\medskip


\noindent   \textbf{Question 8.} Differentiate between any three of the following pairs of families:\pts{$5  \times 3 = 15$}

\begin{parts}
  \item Scrophulariaceae and Meliaceae
  \item Annonaceae and Rutaceae
  \item Moraceae and Liliaceae
  \item Asteraceae and Convolvulaceae
\end{parts}

\bigskip

\begin{center}
   \textbf{SECTION -- C}
\end{center}

\medskip


\noindent   \textbf{Question 9.} Answer the following in not more than 50 words each:\pts{$1  \times 10 = 10$}

\begin{parts}
  \item Differentiate between haplocheilic and syndetocheilic stomata.
  \item What is a feeder and where is it found?
  \item Differentiate between embryogeny of   \textit{Gnetum} and   \textit{Biota}.
  \item Why are Cycadales and Ginkgoales called living fossils?
  \item What is a seed scale complex?
  \item Name a family in which obdiplostemonous condition is found.
  \item Give the botanical name of a medicinal plant of Convolvulaceae.
  \item Name a family in which basal placentation is found.
  \item Write the number of stamens found in a male flower of   \textit{Euphorbia}.
  \item Define hypogynous, epigynous and perigynous flowers.
\end{parts}

\vfill
\rule{\linewidth}{0.4pt}

\begin{center}\small
\href{https://bhu-exam.vercel.app/}{Click Here}\\[4pt]\href{https://me-aryan.vercel.app/}{Click Here}\\[4pt]\href{https://quashan-me.vercel.app/}{Click Here}
\end{center}

\end{document}
